% =====================================================================
%  BÁO CÁO MÔN AI (Trí tuệ Nhân tạo)
%  So sánh Q-learning vs Monte Carlo Control vs Rule-Based System
%  trong bài toán điều khiển tác tử "sói" của mô phỏng sinh thái.
%
%  Biên dịch: pdfLaTeX (đã cấu hình T5 + babel tiếng Việt cho Overleaf).
%  Nếu Overleaf báo lỗi font tiếng Việt: đổi compiler sang XeLaTeX/LuaLaTeX
%  và thay khối inputenc/fontenc bằng:  \usepackage{fontspec}
%
%  Số liệu lấy từ: braincompare_perseed.csv (50 seed x 3000 step),
%  phân tích bằng stats_braincompare.py (ghép cặp) và stats_rliable.py.
% =====================================================================
\documentclass[11pt,a4paper]{article}

\usepackage[utf8]{inputenc}
\usepackage[T5]{fontenc}
\usepackage{lmodern}
\usepackage[vietnamese]{babel}

\usepackage{amsmath,amssymb}
\usepackage{booktabs}
\usepackage{array}
\usepackage{multirow}
\usepackage{graphicx}
\usepackage{geometry}
\usepackage{pgfplots}
\usepackage{caption}
\usepackage{enumitem}
\usepackage[dvipsnames]{xcolor}
\usepackage{hyperref}

\pgfplotsset{compat=1.18}
\geometry{margin=2.4cm}
\hypersetup{colorlinks=true,linkcolor=blue!55!black,citecolor=ForestGreen,urlcolor=blue!55!black}

% \code: in tên file/định danh có dấu gạch dưới mà không cần thoát thủ công.
\newcommand{\code}[1]{\texttt{\detokenize{#1}}}
\newcommand{\catk}{bắt/1000}

\title{\textbf{So sánh Reinforcement Learning với Rule-Based System\\
trong bài toán điều khiển tác tử săn mồi}\\[2mm]
\large Q-learning (TD) \emph{vs} Monte Carlo Control \emph{vs} luật tay (RBS)}

\author{Báo cáo môn Trí tuệ Nhân tạo}
\date{\today}

\begin{document}
\maketitle

\begin{abstract}
\noindent
Báo cáo nghiên cứu bài toán \emph{ra quyết định tuần tự} (sequential decision making) qua một
mô phỏng sinh thái sói--thỏ, trong đó tác tử ``sói'' phải vừa săn mồi vừa tránh chết đói/khát.
Chúng tôi đặt bài toán dưới dạng MDP theo khung \emph{Hierarchical RL / Options}: tác tử học
\emph{chọn macro-action} (1 trong 8 chiến lược luật), mỗi chiến lược tự điều hướng bằng A*. Hai
thuật toán học bảng (tabular) được so sánh trực tiếp --- \textbf{Q-learning} (TD, off-policy,
online) và \textbf{Monte Carlo Control} (return thật, on-policy, offline) --- với cùng không
gian trạng thái / hành động / hàm thưởng; khác biệt duy nhất là \emph{công thức cập nhật}. Cả
hai được đối chứng với baseline \emph{Rule-Based System} (RBS) --- bộ điều khiển chấm điểm bằng
luật tay. Trên \textbf{50 seed
ghép cặp} ($50 \times 3000$ bước), với bộ chỉ số kiểu \emph{rliable} (IQM, Probability of
Improvement) và kiểm định ghép cặp (Wilcoxon signed-rank, bootstrap 95\% CI, Cohen's $d_z$):
khi được huấn luyện \emph{đủ}, \textbf{cả QL lẫn MC đều vượt RBS một cách rõ rệt}
($\text{QL}+51\%$, $P(\text{QL}>\text{RBS})=0.79$; $\text{MC}+35\%$, $P(\text{MC}>\text{RBS})=0.69$),
và \textbf{QL ngang MC} về hiệu quả săn ($P(\text{MC}>\text{QL})=0.41$, khoảng tin cậy chứa $0.5$).
Quan trọng hơn một con số thắng--thua, một \textbf{lưới khử biến $2\times 2$} (thuật toán $\times$ độ
mịn state, kèm trục \emph{ngân sách episode}) tách bạch được ba nguồn khác biệt thật sự: (i) lợi thế
của QL đến từ việc \emph{khai thác được state mịn} (bỏ state mịn $\to$ QL tụt về mức RBS); (ii) MC
\emph{kém hiệu quả mẫu}, cần khoảng $3.75\times$ số episode mới đuổi kịp QL; (iii) MC \emph{nhạy cảm
với reward shaping theo bước} (phạt khát/đói làm \emph{phương sai return} tăng cao $\to$ chết khát
nhiều). Cơ chế \emph{bootstrap} của TD là yếu tố then chốt: nó vừa giúp QL học nhanh từ state mịn,
vừa định vị tín hiệu phạt khát theo từng bước nên quản lý sinh tồn tốt hơn. Kết luận phương pháp:
so sánh ``cùng hàm thưởng tại một ngân sách episode cố định'' dễ gây hiểu nhầm vì trộn lẫn ba yếu
tố trên; phải huấn luyện tới mức đủ (hoặc báo cáo đường cong học) rồi mới so đỉnh.
\end{abstract}

\tableofcontents
\bigskip

% =====================================================================
\section{Giới thiệu}
% =====================================================================

Mục tiêu của báo cáo là trả lời một câu hỏi thực nghiệm: \emph{trong cùng một môi trường và cùng
bộ kỹ năng cơ sở, một tác tử \textbf{học} chính sách (reinforcement learning) có vượt được
một tác tử \textbf{lập trình luật tay} (rule-based) hay không, và giữa hai gia đình thuật toán
học kinh điển --- Temporal Difference (Q-learning) và Monte Carlo --- phương pháp nào phù hợp hơn
cho bài toán này?}

Bài toán nằm trong một mô phỏng sinh thái (lập trình bằng Java/JavaFX). \emph{Về mặt AI}, ta chỉ
cần coi môi trường như một \textbf{Markov Decision Process} (MDP): một bản đồ 2D có nước, cỏ và
con mồi; tác tử sói quan sát cục bộ (lân cận trong tầm nhìn), chọn hành động mỗi bước thời gian,
nhận phần thưởng theo \emph{kết quả} (bắt được mồi, sống sót, chết đói/khát). Phần lập trình môi
trường chỉ được trình bày ở mức cần thiết cho phần AI; trọng tâm là \emph{công thức học}, \emph{thiết
kế hàm phần thưởng} và \emph{phương pháp đảm bảo so sánh công bằng và có ý nghĩa thống kê}.

\paragraph{Ba ``não'' được so sánh.}
\begin{itemize}[leftmargin=1.4em]
  \item \textbf{RBS} (Rule-Based System): baseline không học, chấm điểm các chiến lược ứng viên
        bằng công thức tay rồi chọn điểm cao nhất.
  \item \textbf{QL} (Q-learning): tabular TD control, học \emph{online} mỗi bước.
  \item \textbf{MC} (Monte Carlo Control): học từ return thật, \emph{offline} cuối mỗi episode.
\end{itemize}

\paragraph{Đóng góp.}
(i) Đặt bài toán săn mồi vào khung \emph{Options} để RL và RBS \textbf{dùng chung} tập kỹ năng cơ
sở và thuật toán A*, nhờ đó phép so sánh chỉ phản ánh khác biệt ở \emph{bộ chọn chiến lược}. (ii)
Thiết kế hàm phần thưởng theo kết quả (outcome-based), có \emph{phạt khát/đói} để chống hiện tượng
khai thác lỗ hổng hàm phần thưởng (reward hacking). (iii) Đánh giá theo chuẩn thống kê
hiện đại (rliable IQM \& Probability of Improvement, kiểm định ghép cặp) thay vì chỉ báo trung
bình.

% =====================================================================
\section{Mô hình hoá MDP}\label{sec:mdp}
% =====================================================================

\subsection{Môi trường (mô tả ở mức cần thiết cho phần AI)}

Mỗi ``ván'' (episode) khởi tạo trên bản đồ $576\times 576$ với $50$ cỏ, $10$ thỏ, $3$ sói; bước
thời gian $\Delta t = 0.1$\,s. Sói có nhu cầu \emph{đói} (hunger) và \emph{khát} (thirst) tăng
theo thời gian; nếu một trong hai chạm $100$ thì tác tử chết. Để duy trì sự sống, sói phải bắt mồi
(giảm đói) và uống nước (giảm khát). Đây chính là nguồn \emph{xung đột mục tiêu} khiến bài toán
không tầm thường: thời gian dành cho việc uống là thời gian không thể dành cho việc săn.

\subsection{Ý tưởng then chốt: Hierarchical RL / Options}

Thay vì học điều khiển cấp thấp (gia tốc theo 8 hướng), tác tử học \emph{chọn macro-action} ---
một trong 8 chiến lược luật. Mỗi chiến lược được chọn sẽ tự đảm nhiệm việc điều hướng (qua thuật
toán tìm đường A*) và tự thực hiện các thao tác ăn, uống, sinh sản. Đây là khung \emph{Options}
\cite{sutton1999options}:

\begin{center}
\fbox{\parbox{0.93\textwidth}{\centering
$s$ (trạng thái) $\;\xrightarrow{\;\text{Q-table, }\varepsilon\text{-greedy}\;}\;$
$a \in \{0,\dots,7\}$ $\;\longrightarrow\;$
\code{options[a].updateVelocity(...)} \;(A* + eat/drink/reproduce)}}
\end{center}

\textbf{Hệ quả thiết kế (quan trọng cho tính công bằng):} RL và RBS \emph{chia sẻ} toàn bộ tập
chiến lược và \emph{chia sẻ A*}; chỉ khác \textbf{bộ chọn}. RBS chọn bằng công thức chấm điểm tay;
RL chọn bằng Q-table học được. Do đó so sánh RL vs RBS không lẫn khác biệt về điều hướng --- chỉ
phản ánh chất lượng \emph{chính sách chọn việc}.

\subsection{Không gian hành động (8 macro-action)}

\begin{table}[h]
\centering
\caption{Tám macro-action. Bốn kiểu săn khác nhau ở \emph{lead time} (điểm đón đầu mồi).}
\label{tab:actions}
\small
\begin{tabular}{clp{8.2cm}}
\toprule
Idx & Chiến lược & Ý nghĩa \\
\midrule
0 & \code{Flee}               & chạy trốn kẻ thù \\
1 & \code{SeekWater}          & tìm \& uống nước (có hysteresis) \\
2 & \code{Mate}               & tìm bạn tình để sinh sản \\
3 & \code{Wander}             & lang thang (fallback) \\
4 & \code{Hunter(lead 0.0)}   & \emph{Hunt-direct}: lao thẳng vào mồi \\
5 & \code{Hunter(lead 0.7)}   & \emph{Hunt-short}: đón đầu ngắn ($\equiv$ RBS mặc định) \\
6 & \code{Hunter(lead 1.3)}   & \emph{Hunt-long}: cắt góc, đón xa hơn \\
7 & \code{Hunter(lead 2.0)}   & \emph{Hunt-ambush}: mai phục / đón sâu \\
\bottomrule
\end{tabular}
\end{table}

\textbf{Lý do thiết kế bốn kiểu săn.} RBS chỉ dùng \emph{một} \code{HunterStrategy} cố định (lead
$0.7$). Bốn kiểu săn của RL là \textbf{superset} (tập cha) của RBS, do đó về lý thuyết trần hiệu
năng của RL không thấp hơn trần của RBS: RL có thể chọn kiểu đón đầu phù hợp với từng tình huống
(ở cự ly gần các kiểu hội tụ về nhau, khác biệt chỉ xuất hiện ở cự ly trung bình và xa). Đây chính
là lợi thế cấu trúc mà RL có khả năng khai thác còn RBS thì không.

\subsection{Không gian trạng thái (state encoding)}\label{sec:state}

Quan sát cục bộ được \textbf{rời rạc hoá} thành một chuỗi khoá cho Q-table. Định dạng hiện hành
(phiên bản \emph{rl-struct}):

\begin{center}
\code{P|e3,2|t1,5,1|m4|h1|w0|p0}
\end{center}

\begin{table}[h]
\centering
\caption{Các trường trạng thái của sói/thỏ.}
\label{tab:state}
\small
\begin{tabular}{llp{8.4cm}}
\toprule
Trường & Miền & Ý nghĩa \\
\midrule
role            & \{P,R\}          & P = predator (sói), R = prey (thỏ) \\
enemyDir,Bin    & $0$--$7$/$8$;\,$0$--$3$ & hướng \& khoảng cách kẻ thù gần nhất \\
targetType      & $\{0,1,2\}$      & $0$ không có, $1$ mồi/cỏ, $2$ nước \\
targetDir,Bin   & $0$--$7$/$8$;\,$0$--$3$ & hướng \& khoảng cách mục tiêu \\
preyMoveDir     & $0$--$7$/$8$     & hướng vận tốc mồi (để ``đoán'' mà cắt góc) \\
hungerBin ($h$) & $0$--$3$         & \textbf{4 mức} đói: $<50,\,50\text{--}69,\,70\text{--}89,\,\geq 90$ \\
thirstBin ($w$) & $0$--$2$         & \textbf{3 mức} khát: $<60,\,60\text{--}79,\,\geq 80$ \\
mateBin ($p$)   & $0$--$1$         & đủ điều kiện đẻ \& có bạn tình kề \\
\bottomrule
\end{tabular}
\end{table}

Việc chia \emph{nhiều mức} cho đói (4) và khát (3) là cải tiến của phiên bản \emph{rl-struct}: nó
cho phép chính sách phân biệt được các mức đói vừa, rất đói và sắp chết, từ đó đưa ra quyết định
phù hợp (tiếp tục săn hay tạm dừng để đi uống). Khoảng cách được rời rạc hoá theo tầm nhìn $v$:
$<v/3$ là gần, $<2v/3$ là vừa, còn lại là xa.

\subsection{Hàm thưởng (outcome-based + phạt khát/đói)}\label{sec:reward}

Thưởng theo \textbf{kết quả} của macro-action, \emph{không} dùng shaping khoảng cách dày (shaping
dày phạt oan các kiểu săn trễ như ambush/lead-long vì chúng \emph{tạm thời} lùi để đón). Phần
thưởng mỗi bước (giống hệt nhau ở QL và MC):

\begin{equation}
r =
\underbrace{\begin{cases} -0.02 & \text{(sói: sức ép thời gian)}\\[1mm] +0.05 & \text{(thỏ: sống thêm 1 bước)}\end{cases}}_{\text{nền}}
\;+\;
\underbrace{\begin{cases} +10 & \text{sói bắt được mồi}\\ +1 & \text{thỏ ăn được cỏ}\end{cases}}_{\text{kết quả}}
\;-\; P_{\text{khát}} \;-\; P_{\text{đói}} \;+\; \text{shaping}
\end{equation}

trong đó các thành phần phạt \emph{tuyến tính theo mức}, chỉ kích hoạt khi chỉ số rơi vào vùng nguy
hiểm (vì vậy tác tử không thể lợi dụng để tích lũy phần thưởng):
\begin{align}
P_{\text{khát}} &= 0.2\cdot\tfrac{(\text{thirst}-60)^{+}}{40}
                 + 0.4\cdot\tfrac{(\text{thirst}-80)^{+}}{20}, &
P_{\text{đói}}  &= 0.2\cdot\tfrac{(\text{hunger}-50)^{+}}{50},
\end{align}
với $(x)^{+}=\max(x,0)$. Shaping tiến-gần mục tiêu là \emph{bất đối xứng} ($+0.03\cdot\text{closed}$
chỉ khi lại gần, lùi thì $0$ chứ không phạt); thỏ được thêm $+0.06\cdot\Delta$ khoảng cách với
sói. Khi chết: phần thưởng cuối $-10$ (terminal).

\begin{table}[h]
\centering
\caption{Tổng hợp reward. Nguyên tắc thang bậc: sự kiện mục tiêu ($\pm 10$) $\gg$ shaping ($\sim 0.03$--$0.06$/bước).}
\label{tab:reward}
\small
\begin{tabular}{lcc}
\toprule
Sự kiện & Sói (PREDATOR) & Thỏ (PREY) \\
\midrule
Mỗi bước (nền)        & $-0.02$ & $+0.05$ \\
Bắt được mồi          & $+10$ & --- \\
Ăn cỏ                 & --- & $+1$ \\
Phạt khát (vùng nguy) & tới $\approx -0.6$/bước & tới $\approx -0.6$/bước \\
Phạt đói (chống camp) & tới $\approx -0.2$/bước & --- \\
Lại gần mục tiêu      & $+0.03\cdot$closed & $+0.03\cdot$closed \\
Tăng cách với sói     & --- & $+0.06\cdot\Delta$ \\
Chết (terminal)       & $-10$ & $-10$ \\
\bottomrule
\end{tabular}
\end{table}

\paragraph{Lý do dùng phạt khát/đói thay cho thưởng uống.}
Phiên bản trước cộng $+4$ cho mỗi lần uống, nhưng cơ chế này tạo ra hiện tượng khai thác lỗ hổng
hàm phần thưởng (\emph{reward hacking}): tác tử học cách trú lại cạnh nguồn nước để liên tục thu
phần thưởng uống và phần thưởng sống sót, thay vì đi săn. Thiết kế hiện tại chỉ \textbf{phạt} khát
trong vùng nguy hiểm và bổ sung \textbf{phạt đói đối xứng}, khiến chiến lược trú cạnh nước không
còn có lợi và buộc tác tử phải đi săn để giảm đói. Đây là một lưu ý quan trọng về thiết kế hàm
phần thưởng, đồng thời là một yếu tố gây nhiễu (\emph{confound}) sẽ được phân tích ở phần thảo
luận, do Monte Carlo phản ứng với tín hiệu phạt khát theo cách rất khác Q-learning.

% =====================================================================
\section{Hai thuật toán học: TD vs Monte Carlo}\label{sec:algos}
% =====================================================================

Cả hai đều cập nhật dạng $Q(s,a) \leftarrow Q(s,a) + \alpha\,[\,\text{target} - Q(s,a)\,]$.
\textbf{Khác biệt duy nhất} là \emph{target}.

\subsection{Q-learning (TD, bootstrap) --- cập nhật MỖI bước}
\begin{equation}
Q(s,a)\;\leftarrow\;Q(s,a) + \alpha\Big[\,\underbrace{r + \gamma\max_{a'}Q(s',a')}_{\text{target (bootstrap)}} - Q(s,a)\Big]
\end{equation}
Phương pháp này sử dụng chính ước lượng hiện có $Q(s')$ làm xấp xỉ cho ``phần thưởng tương lai''.
Tính chất: \emph{off-policy} (dùng toán tử $\max$), \emph{online} (cập nhật ngay từng bước),
\emph{phương sai thấp nhưng có độ chệch (bias)}. Tại trạng thái cuối (tác tử chết): target $=r$
(không bootstrap).

\subsection{Monte Carlo Control (return thật) --- cập nhật CUỐI episode}
\begin{equation}
G_t = \sum_{k=0}^{T-t-1}\gamma^{k}\,r_{t+k+1}, \qquad
Q(s_t,a_t)\;\leftarrow\;Q(s_t,a_t) + \alpha\,[\,G_t - Q(s_t,a_t)\,]
\end{equation}
Dùng \emph{return thật} $G_t$ (không bootstrap), \emph{on-policy}, \emph{offline} (đợi episode kết
thúc). Cần lưu lại quỹ đạo (trajectory) các cặp $(s,a)$; phần thưởng được điền sau theo cơ chế trễ
một bước. Tính chất: \emph{không chệch (unbiased) nhưng phương sai cao}.

\begin{table}[h]
\centering
\caption{Bản chất hai thuật toán --- giải thích kỳ vọng về kết quả.}
\label{tab:tdvsmc}
\small
\begin{tabular}{lll}
\toprule
& \textbf{Q-learning (TD)} & \textbf{Monte Carlo} \\
\midrule
Target                 & $r+\gamma\max_{a'}Q(s',a')$ & return thật $G_t$ \\
Bootstrap              & Có & Không \\
Bias / Variance        & có bias / \emph{variance thấp} & không bias / \emph{variance cao} \\
Policy                 & off-policy & on-policy \\
Thời điểm học          & online, mỗi bước & offline, cuối episode \\
Credit qua horizon dài & dần dần qua nhiều episode & toàn bộ episode trong 1 lượt \\
\bottomrule
\end{tabular}
\end{table}

\subsection{Tham số học \& pipeline huấn luyện}\label{sec:train}

\begin{itemize}[leftmargin=1.4em,itemsep=1pt]
  \item $\alpha=0.2$ (QL) / $0.1$ (MC);\; $\gamma=0.99$ (horizon hiệu dụng $\approx 1/(1-\gamma)
        \approx 100$ bước --- chọn để phủ được thang thời gian đói/khát kéo dài hàng trăm bước).
  \item $\varepsilon$ giảm tuyến tính $1.0\!\to\!0.05$, chạm sàn ở $\sim 80\%$ số episode;
        \textbf{eval $\varepsilon=0$} (tất định, để công bằng với RBS).
  \item \textbf{Giao thức huấn luyện thống nhất (best-response một pha).} Vì trong mọi phép đánh
        giá con mồi \emph{luôn là RBS}, ta huấn luyện sói theo \emph{một pha duy nhất}: sói học
        chống thỏ RBS cố định. Cách này loại bỏ tính phi-dừng (non-stationary) của lối huấn luyện
        xen kẽ trước đây, đồng thời cho phép \textbf{dùng cùng độ dài episode cho cả QL và MC}.
  \item \textbf{Đã khử confound độ dài episode.} Mọi cấu hình trong báo cáo này train với
        \emph{cùng} $800$ bước/episode (phiên bản trước dùng MC$=400$, QL$=600$); nhờ vậy khác biệt
        kết quả không còn lẫn biến độ-dài-episode. Ngân sách mặc định $800$ episode; riêng MC được
        khảo sát thêm ở $3000$ episode (xem lưới khử biến, \S\ref{sec:grid}).
\end{itemize}

% =====================================================================
\section{Baseline RBS (Rule-Based System)}
% =====================================================================

RBS không học: mỗi bước, sói chấm điểm mọi chiến lược ứng viên rồi chọn điểm cao nhất (có
\emph{commit bonus} giảm dần theo thời gian để chống nhảy qua lại / flicker). Bảng chấm điểm của
sói (rút gọn):

\begin{table}[h]
\centering
\caption{Logic chấm điểm RBS cho sói. Hunter của RBS \emph{cố định lead $0.7$} (chỉ 1 kiểu săn).}
\label{tab:rbs}
\small
\begin{tabular}{lp{9.6cm}}
\toprule
Chiến lược & Điều kiện \& điểm \\
\midrule
Flee      & có đe doạ $\Rightarrow 100$ (luôn thắng) \\
SeekWater & emergency nếu thirst $\geq$ critical; commit nếu đang uống dở; seek theo thirst \\
Hunter    & chưa đói ($<50$) $\Rightarrow 0$; không thấy mồi $\Rightarrow 0$; mồi rất gần $\Rightarrow$ tấn công áp sát; còn lại $\max(0.75,\text{hunger}/100)+0.5\cdot$closeness \\
Mate      & đủ điều kiện $\Rightarrow 0.45$ \\
Wander    & $0.3$ (nền, fallback) \\
\bottomrule
\end{tabular}
\end{table}

Đây chính là cơ sở để so sánh công bằng: RBS và RL chỉ khác \textbf{bộ chọn} (công thức tay vs
Q-table), mọi thứ khác (tập chiến lược, A*, ngưỡng sinh thái) giữ nguyên.

% =====================================================================
\section{Thiết kế thí nghiệm so sánh}\label{sec:exp}
% =====================================================================

\subsection{Nguyên tắc: cùng điều kiện, so theo cặp}

So \emph{ba bộ điều khiển sói} trên CÙNG điều kiện, \textbf{ghép cặp theo seed}: cùng seed $\to$
cùng layout thế giới (vị trí spawn, cỏ, nước) cho cả ba não. Biến \emph{duy nhất} thay đổi là bộ
điều khiển của sói; con mồi (thỏ) \emph{luôn RBS} ở mọi nhánh.

\begin{itemize}[leftmargin=1.4em,itemsep=1pt]
  \item \textbf{Thế giới cô lập} (chỉ gồm sói, thỏ và cỏ): số lần bắt mồi (catch) đếm được phản ánh
        chính xác số mồi do sói bắt được.
  \item Sói thả \emph{đói sẵn} ($\text{hunger}=55$, ngưỡng săn $50$) để cả 3 não khởi động cùng điều
        kiện; \textbf{bù thỏ về sàn $10$} mỗi bước (áp lực mồi không tắt), \textbf{không bù sói} (để
        đo tuổi thọ \& chết đói/khát).
  \item RNG của $\varepsilon$ lấy seed riêng tất định (theo seed $\times$ hằng $+$ chỉ số sói), không
        lấy từ RNG thế giới --- nhờ vậy vị trí spawn giống hệt giữa 3 nhánh.
  \item Eval $\varepsilon=0$: RL tất định như RBS (bỏ handicap bước ngẫu nhiên).
\end{itemize}

\subsection{Cấu hình chạy}
$50$ seed $\times\,3000$ bước ($\approx 300$\,s mô phỏng/seed), công cụ headless \code{BrainCompare}.
Reset bộ đếm A* ngay trước vòng mô phỏng (spawn không tính vào A*).

\subsection{Các metric}
\begin{table}[h]
\centering
\caption{Bộ metric. \emph{\catk\ step} là chỉ số chính.}
\label{tab:metrics}
\small
\begin{tabular}{lcl}
\toprule
Metric & Hướng tốt & Ý nghĩa \\
\midrule
\catk\ step             & cao  & hiệu quả săn $=$ catches$\cdot 1000/$steps \;(\textbf{chính}) \\
Tuổi thọ TB sói (step)  & cao  & trung bình (deathStep $-$ birthStep) \\
Sói chết đói            & thấp & chết do hunger \\
Sói chết khát           & thấp & chết do thirst \\
Step tới lần bắt đầu     & thấp & $=$ maxStep nghĩa là \emph{không} bắt được con nào \\
CPU ms/step             & thấp & chi phí tính toán (đại diện) \\
A* gọi/step, node/step  & (tham khảo) & tải tìm đường --- 3 não chung A* nên so công bằng \\
\bottomrule
\end{tabular}
\end{table}

\subsection{Phương pháp thống kê}\label{sec:stats}

Thiết kế là \textbf{matched-pairs} (cùng seed $\to$ cùng layout), nên không được dùng kiểm định
không-ghép-cặp. Chúng tôi báo cáo hai tầng bổ sung nhau:

\paragraph{(A) Kiểm định ghép cặp \& độ lớn hiệu ứng} (\code{stats_braincompare.py}).
Với hiệu $d_i = A_i - B_i$ trên từng seed: \emph{Wilcoxon signed-rank} (phi tham số, xấp xỉ chuẩn
có hiệu chỉnh tie + liên tục) là kiểm định \emph{chính}; \emph{paired $t$} (xấp xỉ chuẩn) để đối
chiếu; \emph{Cohen's $d_z$} cho độ lớn hiệu ứng; \emph{bootstrap $95\%$ CI} ($B=20000$) cho hiệu
trung bình \& trung vị; và đếm thắng/hòa/thua.

\paragraph{(B) Đánh giá kiểu rliable} \cite{agarwal2021rliable} (\code{stats_rliable.py}).
Báo cáo \emph{Median / IQM / Mean} kèm $95\%$ stratified-bootstrap CI cho từng não. \textbf{IQM}
(Interquartile Mean --- trung bình $50\%$ giữa) là chỉ số được rliable khuyến nghị vì \emph{bền
vững với outlier} hơn mean và \emph{nhạy} hơn median. Ngoài ra báo cáo \emph{Probability of
Improvement} $P(A>B)$ kèm CI, và win-rate ghép cặp khai thác cùng-seed.

% =====================================================================
\section{Kết quả}\label{sec:results}
% =====================================================================

\subsection{Metric chính --- best-vs-best (trung bình $\pm$ độ lệch chuẩn trên 50 seed)}

Bảng~\ref{tab:main} so ba não ở cấu hình \emph{tốt nhất / huấn luyện đủ} của mỗi bên, trong
\emph{cùng một lần chạy ghép cặp} (cùng 50 seed): RBS, QL với state mịn ($800$ ep), và MC với state
mịn nhưng được huấn luyện đủ ($3000$ ep --- xem lý do ở \S\ref{sec:grid}).

\begin{table}[h]
\centering
\caption{Bảng metric tổng hợp (RBS \emph{vs} QL@mịn \emph{vs} MC@mịn-$3000$ep). \textbf{In đậm} = tốt nhất theo hàng.}
\label{tab:main}
\small
\begin{tabular}{lccc}
\toprule
Metric & RBS & QL & MC \\
\midrule
\catk\ step (cao=tốt)   & $4.65\pm1.98$ & $\mathbf{7.01\pm2.37}$ & $6.27\pm2.81$ \\
Tuổi thọ TB sói (step)  & $\mathbf{2722\pm313}$ & $2530\pm454$ & $2125\pm464$ \\
Sói chết đói            & $\mathbf{0.32\pm0.55}$ & $0.38\pm0.53$ & $0.66\pm0.75$ \\
Sói chết khát           & $\mathbf{0.40\pm0.64}$ & $0.66\pm0.82$ & $1.26\pm0.90$ \\
Step tới lần bắt đầu     & $\mathbf{69\pm89}$ & $117\pm110$ & $130\pm151$ \\
CPU ms/step (thấp=tốt)  & $4.34\pm1.40$ & $2.82\pm1.26$ & $\mathbf{2.49\pm0.95}$ \\
A* gọi/step             & $4.03\pm1.40$ & $2.56\pm1.35$ & $\mathbf{2.17\pm0.87}$ \\
A* node/step            & $1244.8\pm424.7$ & $785.9\pm380.7$ & $\mathbf{680.2\pm293.6}$ \\
\bottomrule
\end{tabular}
\end{table}

\noindent
Cả hai não học đều bắt mồi hiệu quả hơn RBS rõ rệt, và \textbf{đều rẻ hơn RBS về tính toán}: chúng
gọi A* ít hơn do chính sách ổn định, ít đổi mục tiêu thừa. Cái giá của RL là \emph{sinh tồn}: cả QL
lẫn MC chết khát nhiều hơn RBS (RBS là kẻ sống sót thận trọng nhưng săn kém). Đáng chú ý, MC chết
khát gần gấp đôi QL ($1.26$ so với $0.66$) và tuổi thọ thấp nhất --- một dấu hiệu sẽ được mổ xẻ ở
\S\ref{sec:disc}.

\subsection{Tổng hợp kiểu rliable (metric: \catk\ step)}

\begin{table}[h]
\centering
\caption{Median / IQM / Mean kèm $95\%$ bootstrap CI. IQM là chỉ số khuyến nghị.}
\label{tab:rliable}
\small
\begin{tabular}{lccc}
\toprule
Não & Median [95\% CI] & \textbf{IQM} [95\% CI] & Mean [95\% CI] \\
\midrule
RBS & $4.65\,[4.12,5.21]$ & $\mathbf{4.36\,[3.79,5.03]}$ & $4.65\,[4.11,5.20]$ \\
QL  & $7.01\,[6.36,7.66]$ & $\mathbf{6.96\,[6.27,7.68]}$ & $7.01\,[6.36,7.65]$ \\
MC  & $6.27\,[5.51,7.03]$ & $\mathbf{6.19\,[5.45,6.96]}$ & $6.27\,[5.53,7.05]$ \\
\bottomrule
\end{tabular}
\end{table}

\begin{figure}[h]
\centering
\begin{tikzpicture}
\begin{axis}[
  width=0.74\textwidth, height=6.2cm,
  ybar, bar width=26pt,
  symbolic x coords={RBS,QL,MC},
  xtick=data,
  ymin=0, ymax=9,
  ylabel={IQM \catk\ step},
  nodes near coords, nodes near coords align={vertical},
  every node near coord/.append style={font=\small},
  enlarge x limits=0.32,
  error bars/y dir=both, error bars/y explicit,
]
\addplot+[draw=black,fill=gray!35] coordinates {
  (RBS,4.36) +- (0,0.62)
  (QL,6.96)  +- (0,0.70)
  (MC,6.19)  +- (0,0.75)
};
\end{axis}
\end{tikzpicture}
\caption{IQM hiệu quả săn với thanh sai số $95\%$ CI (bootstrap, $\approx$ đối xứng), cấu hình
best-vs-best. CI của \emph{cả QL lẫn MC} tách khỏi RBS; CI của QL và MC \emph{chồng nhau} --- hai
phương pháp học hòa nhau khi MC được huấn luyện đủ.}
\label{fig:iqm}
\end{figure}

\subsection{Kiểm định ghép cặp}

\begin{table}[h]
\centering
\caption{So sánh ghép cặp trên \catk\ step ($n=50$). CI bootstrap $95\%$ cho hiệu trung bình.}
\label{tab:paired}
\small
\begin{tabular}{lcccc}
\toprule
Cặp (A$-$B) & Hiệu TB [95\% CI] & Wilcoxon $z$ ($p$) & $d_z$ & Thắng/Hòa/Thua \\
\midrule
QL $-$ RBS & $+2.36\,[+1.47,+3.23]$ & $+4.30\;(1.7\!\times\!10^{-5})$ & $+0.73$ & $41/1/8$ \\
MC $-$ RBS & $+1.62\,[+0.70,+2.58]$ & $+2.98\;(2.9\!\times\!10^{-3})$ & $+0.47$ & $34/2/14$ \\
MC $-$ QL  & $-0.74\,[-1.73,+0.32]$ & $-1.76\;(7.9\!\times\!10^{-2})$ & $-0.20$ & $18/2/30$ \\
\bottomrule
\end{tabular}
\end{table}

\begin{itemize}[leftmargin=1.4em,itemsep=1pt]
  \item \textbf{QL $>$ RBS}: hiệu $+2.36$ \catk\ ($+51\%$), CI \emph{không chứa $0$}, $d_z=0.73$
        (hiệu ứng \emph{lớn}), thắng $41/50$ seed. Rất có ý nghĩa.
  \item \textbf{MC $>$ RBS}: hiệu $+1.62$ ($+35\%$), CI \emph{không chứa $0$}, Wilcoxon $p=0.003$,
        $d_z=0.47$ (hiệu ứng \emph{vừa}), thắng $34/50$. Có ý nghĩa --- \emph{khác hẳn} kết quả
        thiếu-train (xem \S\ref{sec:grid}).
  \item \textbf{MC vs QL}: hiệu $-0.74$, \textbf{CI chứa $0$}, Wilcoxon $p=0.079$, $d_z=-0.20$ (nhỏ),
        $18/2/30$ $\Rightarrow$ \emph{không} khác biệt có ý nghĩa: QL và MC \textbf{hòa} về hiệu quả săn.
\end{itemize}

\subsection{Probability of Improvement (rliable)}

\begin{table}[h]
\centering
\caption{$P(A>B)$ kèm $95\%$ CI (Mann-Whitney) và win-rate ghép cặp. $0.5 =$ không khác.}
\label{tab:poi}
\small
\begin{tabular}{lccc}
\toprule
Cặp & $P(A>B)$ [95\% CI] & win-rate ghép cặp & Kết luận \\
\midrule
$P(\text{QL}>\text{RBS})$ & $0.789\,[0.694,0.873]$ & $0.830$ & QL tốt hơn \\
$P(\text{MC}>\text{RBS})$ & $0.688\,[0.581,0.789]$ & $0.700$ & MC tốt hơn \\
$P(\text{MC}>\text{QL})$  & $0.411\,[0.304,0.525]$ & $0.380$ & $\approx$ ngang (CI chứa $0.5$) \\
\bottomrule
\end{tabular}
\end{table}

Cả hai tầng phương pháp \textbf{nhất quán}: $\text{QL} \approx \text{MC} \;>\; \text{RBS}$ ---
cả hai thuật toán học đều vượt baseline một cách có ý nghĩa, và hòa nhau ở chỉ số chính (khi MC
được huấn luyện đủ).

\subsection{Lưới khử biến: vì sao QL thắng và MC ``hụt hơi''?}\label{sec:grid}

Con số best-vs-best ở trên (QL $\approx$ MC $>$ RBS) chưa giải thích \emph{tại sao}, và đặc biệt vì
sao MC cần tới $3000$ episode trong khi QL chỉ cần $800$. Để tách bạch, chúng tôi chạy một
\textbf{lưới khử biến} (ablation), giữ \emph{cố định} reward + $\gamma$ + giao thức ($800\times 800$,
một pha, thỏ RBS) và chỉ thay đổi hai trục: \textbf{thuật toán} (QL/MC) $\times$ \textbf{độ mịn của
state} (thô $1$-bit \emph{vs} mịn rl-struct), cộng trục \emph{ngân sách episode} cho MC. Mỗi ô là một
lần train-từ-rỗng rồi eval $50$ seed; IQM \catk\ step:

\begin{table}[h]
\centering
\caption{Lưới khử biến (IQM \catk\ step, $50$ seed). Mọi ô: một pha, $800$ bước/ep, reward+$\gamma$
rl-struct. RBS $\approx 4.4$ làm mốc. Tham chiếu (có confound): MC@thô với \emph{reward cũ} +
pipeline xen kẽ đạt $7.14$.}
\label{tab:grid}
\small
\begin{tabular}{lcc}
\toprule
 & state THÔ ($1$-bit) & state MỊN (rl-struct) \\
\midrule
QL ($800$ ep)  & $4.21$ \;($\approx$ RBS) & $\mathbf{6.96}$ \\
MC ($800$ ep)  & $3.95$ & $3.62$ \\
MC ($3000$ ep) & --- & $6.19$ \\
\bottomrule
\end{tabular}
\end{table}

\noindent Lưới này tách được \textbf{ba nguồn khác biệt} (trước đây bị trộn lẫn):

\begin{enumerate}[leftmargin=1.6em,itemsep=2pt]
  \item \textbf{Lợi thế của QL = khai thác được state mịn.} Giữ nguyên reward, chỉ thô-hoá state,
        QL tụt từ $6.96 \to 4.21$ (ngang RBS, $P(\text{QL}>\text{RBS})=0.44$). Tức năng lực vượt
        trội của QL \emph{không} đến từ hàm thưởng mà từ việc \emph{phân giải trạng thái mịn hơn}
        --- điều mà cập nhật TD (bootstrap) khai thác được ngay cả khi mỗi trạng thái ít mẫu.
  \item \textbf{MC kém hiệu quả mẫu.} Với cùng $800$ ep, MC@mịn chỉ đạt $3.62$; phải tăng lên
        $3000$ ep (khoảng $3.75\times$) MC mới hồi về $6.19$ và hòa QL. MC \emph{không} hỏng về bản
        chất --- nó chỉ cần nhiều mẫu hơn hẳn để trung bình hoá return.
  \item \textbf{Thủ phạm là phương sai return, không phải kích thước state.} Một giả thuyết tự
        nhiên là ``state mịn quá lớn nên MC thiếu mẫu''. Nhưng MC@thô (bảng chỉ $286$ trạng thái,
        $800$ ep \emph{thừa} mẫu) vẫn chỉ đạt $3.95 \approx$ MC@mịn $3.62$. Bảng nhỏ, đủ mẫu, mà vẫn
        kẹt $\Rightarrow$ vấn đề \emph{không} phải số trạng thái. Nguyên nhân là \textbf{phương sai
        cao của return} sinh ra bởi các khoản \emph{phạt khát/đói theo bước}: chúng nhỏ, rải đều,
        xuất hiện muộn, làm $G_t$ nhiễu mạnh; MC (trung bình return cả episode) cần rất nhiều mẫu
        mới lọc nổi. Bằng chứng: MC \emph{chết khát cao ở mọi ô} dùng reward này --- đúng hành vi mà
        khoản phạt định dạy thì MC lại không học được. (Đối chiếu: với \emph{reward cũ} gọn hơn,
        ít phạt theo bước, MC@thô từng đạt $7.14$ --- nhất quán với giả thuyết phương sai, dù ô đó
        còn lẫn biến giao thức nên chỉ mang tính tham chiếu.)
\end{enumerate}

% =====================================================================
\section{Thảo luận}\label{sec:disc}
% =====================================================================

\subsection{Phân tích: nguyên nhân QL vượt RBS}
QL khai thác được toàn bộ \emph{superset} gồm bốn kiểu săn: chính sách học được chọn lead time phù
hợp theo cự ly và hướng di chuyển của mồi --- điều mà RBS, vốn cố định một kiểu săn, không thực
hiện được. Đáng chú ý, QL đạt hiệu quả cao hơn trong khi số lần gọi A* lại thấp hơn khoảng $40\%$
(Bảng~\ref{tab:main}), cho thấy chính sách ra quyết định ổn định và ít dao động giữa các chiến
lược. Tuổi thọ của QL thấp hơn RBS một chút (chết khát $0.66$ so với $0.40$): QL chấp nhận rủi ro
về khát cao hơn để đổi lấy nhiều cơ hội săn hơn --- một sự đánh đổi hợp lý, vì chỉ số mục tiêu của
bài toán là hiệu quả săn. Nhờ phân giải state mịn (\S\ref{sec:grid}), QL chọn được kiểu đón đầu theo
từng tình huống, đó là phần năng lực mà RBS (một kiểu săn cố định) không có.

\subsection{Phân tích: hồ sơ sinh tồn của MC và khoảng cách hiệu quả mẫu}
Khi được huấn luyện đủ, MC \emph{hòa} QL về hiệu quả săn (Bảng~\ref{tab:main}), nhưng vẫn lộ một
điểm yếu nhất quán: \textbf{MC chết khát nhiều hơn} ($1.26$ so với $0.66$ của QL; và lên tới $1.92$
khi thiếu train) kéo tuổi thọ xuống thấp nhất ($2125$ bước). Quan trọng hơn, lưới khử biến
(\S\ref{sec:grid}) cho thấy ở \emph{cùng} ngân sách $800$ ep, MC chỉ đạt $\approx 3.6$--$4.0$, kém
QL hẳn một bậc. Ba quan sát có liên hệ giải thích điều này:

\begin{enumerate}[leftmargin=1.6em,itemsep=2pt]
  \item \textbf{Phương sai cao của ước lượng return --- thủ phạm chính.} MC cập nhật $Q(s,a)$ theo
        $G_t$, tổng chiết khấu của \emph{toàn bộ} phần còn lại episode. Các khoản phạt khát/đói nhỏ
        ($\sim 0.2$--$0.6$/bước), rải đều và đến muộn, cộng với cú $-10$ khi chết, làm $G_t$ \emph{nhiễu
        rất mạnh}; trạng thái ``khát cao'' do đó khó có ước lượng $Q$ ổn định để học được quy tắc
        ``cần đi uống''. Lưới khử biến chứng minh đây là nguyên nhân \emph{thực sự}, không phải kích
        thước không gian trạng thái: ngay cả với bảng thô chỉ $286$ trạng thái (thừa mẫu ở $800$ ep),
        MC vẫn kẹt ở $3.95$. Ngược lại, TD lan truyền phạt khát \emph{cục bộ theo từng bước} qua
        bootstrap, phương sai thấp, nên QL học quản lý khát tốt hơn hẳn.
  \item \textbf{Kém hiệu quả mẫu, nhưng không hỏng về bản chất.} Vì cần trung bình hoá phương sai
        nói trên, MC đòi hỏi nhiều mẫu hơn nhiều: tăng từ $800$ lên $3000$ ep ($\approx 3.75\times$)
        đưa MC@mịn từ $3.62$ lên $6.19$, đủ hòa QL. Đây là biểu hiện kinh điển của khoảng cách hiệu
        quả mẫu giữa Monte Carlo (phương sai cao) và TD.
  \item \textbf{Đã khử confound độ dài episode.} Phiên bản trước train MC với $400$ bước, QL với
        $600$ bước --- một biến gây nhiễu. Trong báo cáo này, \emph{mọi} cấu hình dùng chung $800$
        bước/ep (\S\ref{sec:train}), nên khác biệt còn lại phản ánh đúng bản chất thuật toán chứ
        không phải độ-dài-episode. Câu hỏi ``under-training hay bản chất'' từng để ngỏ nay đã được
        trả lời bằng trục ngân sách episode của lưới: \emph{cả hai} --- MC kém ở ngân sách thấp
        (sample-inefficiency) nhưng hồi phục khi đủ mẫu, và luôn yếu hơn ở quản lý khát (phương sai).
\end{enumerate}

\paragraph{Tương tác giữa thuật toán và thiết kế hàm phần thưởng.} Việc chuyển từ \emph{thưởng uống
$+4$} sang \emph{phạt khát} (\S\ref{sec:reward}) đã ngăn được hiện tượng khai thác lỗ hổng ``trú
cạnh nguồn nước'', nhưng đồng thời biến mục tiêu ``tránh chết khát'' thành một tín hiệu phạt có giá
trị nhỏ, phân bố đều và xuất hiện muộn --- đúng dạng tín hiệu mà MC xử lý kém nhất. Nói cách khác,
cùng một hàm phần thưởng có thể phù hợp với TD nhưng bất lợi cho MC. Lưới khử biến củng cố trực tiếp
điểm này: ở \emph{cùng} state thô, MC dùng reward shaping rl-struct chỉ đạt $3.95$, trong khi với
reward cũ (ít phạt theo bước) từng đạt $7.14$. Đây là một minh hoạ cụ thể cho nguyên lý: việc lựa
chọn thuật toán và thiết kế hàm phần thưởng \emph{không độc lập với nhau}.

% =====================================================================
\section{Hạn chế \& hướng phát triển}
% =====================================================================
\begin{itemize}[leftmargin=1.4em,itemsep=1pt]
  \item \textbf{Giảm phương sai cho MC:} đã xác định phương sai return là thủ phạm chính; bước tiếp
        theo là dùng baseline / return chuẩn hoá, hoặc chuyển sang $n$-step / TD($\lambda$) để nối
        nhịp TD--MC --- kỳ vọng thu hẹp khoảng cách hiệu quả mẫu.
  \item \textbf{Đóng nốt lưới:} ô MC@thô ở $3000$ ep (chưa chạy) và ô MC@thô với \emph{reward cũ}
        ở giao thức một-pha (khử confound còn lại của con số tham chiếu $7.14$).
  \item \textbf{Đường cong học:} báo cáo catches/episode theo thời gian train để định lượng tốc độ
        hội tụ TD \emph{vs} MC (lưới hiện chỉ chụp ở $800$ và $3000$ ep).
  \item \textbf{Đo trần môi trường (oracle):} có trần mới tính được \emph{optimality gap} chuẩn hoá
        $[0,1]$ của rliable (hiện bỏ qua vì chưa đo trần).
\end{itemize}

% =====================================================================
\section{Kết luận}
% =====================================================================
Trên cùng môi trường, cùng không gian trạng thái, hành động, hàm phần thưởng và cùng thuật toán A*,
\textbf{cả hai thuật toán học đều vượt baseline luật tay một cách rõ rệt} khi được huấn luyện đủ
($\text{QL}+51\%$, $P(\text{QL}>\text{RBS})=0.79$, $d_z=0.73$; $\text{MC}+35\%$,
$P(\text{MC}>\text{RBS})=0.69$), đồng thời rẻ hơn RBS về tính toán; và \textbf{QL ngang MC} ở hiệu
quả săn ($P(\text{MC}>\text{QL})=0.41$). Tuy nhiên, lưới khử biến cho thấy ``ngang nhau ở đỉnh''
che giấu ba khác biệt cốt lõi: (i) lợi thế của QL đến từ khả năng \emph{khai thác state mịn} (bỏ đi
thì QL tụt về mức RBS); (ii) MC \emph{kém hiệu quả mẫu}, cần khoảng $3.75\times$ số episode mới đuổi
kịp; (iii) MC \emph{nhạy cảm với reward shaping theo bước} do phương sai return cao, biểu hiện ở việc
chết khát nhiều --- đúng thứ mà cập nhật bootstrap của TD xử lý tốt. Ba kết luận rút ra: (1) RL có
thể vượt hệ luật tay khi được cấp đúng lợi thế về không gian hành động (superset) và phân giải trạng
thái; (2) giữa các thuật toán RL, sự phù hợp với \emph{cấu trúc hàm thưởng} và \emph{ngân sách mẫu}
quan trọng hơn nhãn ``có phải RL hay không''; và (3) về phương pháp luận, so sánh hai thuật toán ở
\emph{một} ngân sách episode cố định dễ gây hiểu nhầm --- phải huấn luyện tới mức đủ hoặc báo cáo
đường cong học thì kết luận mới vững.

\paragraph{Tái lập.} Số liệu sinh từ \code{BrainCompare} ($50\times3000$), lưu
\code{braincompare_perseed.csv}; phân tích bằng \code{stats_braincompare.py} (ghép cặp) và
\code{stats_rliable.py} (rliable). Eval $\varepsilon=0$, seed ghép cặp.

% =====================================================================
\begin{thebibliography}{9}
\bibitem{sutton2018rl} R.~S. Sutton and A.~G. Barto.
  \emph{Reinforcement Learning: An Introduction}, 2nd ed. MIT Press, 2018.
\bibitem{watkins1992q} C.~J.~C.~H. Watkins and P. Dayan.
  Q-learning. \emph{Machine Learning}, 8(3--4):279--292, 1992.
\bibitem{sutton1999options} R.~S. Sutton, D. Precup, and S. Singh.
  Between MDPs and semi-MDPs: A framework for temporal abstraction in reinforcement learning.
  \emph{Artificial Intelligence}, 112(1--2):181--211, 1999.
\bibitem{agarwal2021rliable} R. Agarwal, M. Schwarzer, P.~S. Castro, A. Courville, and M.~G. Bellemare.
  Deep Reinforcement Learning at the Edge of the Statistical Precipice.
  \emph{NeurIPS}, 2021.
\bibitem{wilcoxon1945} F. Wilcoxon.
  Individual comparisons by ranking methods. \emph{Biometrics Bulletin}, 1(6):80--83, 1945.
\bibitem{efron1993boot} B. Efron and R.~J. Tibshirani.
  \emph{An Introduction to the Bootstrap}. Chapman \& Hall, 1993.
\end{thebibliography}

\end{document}
